\section{Related Work}
\label{sec:related}

\paragraph{Prompt order and position effects.} That the position of information
in a prompt changes what a model does is established in text. Few-shot example
order alone swings accuracy across a wide range \citep{fantasticallyordered2022},
and long-context models recover facts placed at the ends of a context better
than facts placed in the middle \citep{lostinthemiddle2024}, and both concern
where evidence sits relative to the query. The question-first
paradox is the multimodal counterpart with a sharper edge: our two layouts hold
the same tokens in the same quantity, so context length, truncation and
retrieval difficulty cannot distinguish them. Re-reading the question also helps
text reasoning \citep{rereading2024}, the same intervention as our question echo
reached from a different direction; that work motivates it as deliberation,
whereas our diagnosis says it restores an attention path.

\paragraph{Prompt optimization.} An alternative to reasoning about layout is to
search over prompt text. DSPy compiles declarative pipelines into tuned prompts
\citep{dspy2024}, and GEPA evolves prompts by reflecting on failures
\citep{gepa2025}. We use GEPA not as a competitor but as a test of the
diagnosis, because it can rewrite the system message and cannot move the image:
Section~\ref{sec:gepa} shows it recovers part of the loss under the broken
ordering and nothing under the working one, which is what a read-out account
predicts and a wording account does not.

\paragraph{Interpretability of vision-language models.} We read intermediate
states with a logit lens \citep{logitlens2020,tunedlens2023} and establish
causality by severing attention edges, the method used to localize factual
recall in language models \citep{attnknockout2023,rome2022}. Our use is
diagnostic rather than editing: we do not intervene to change behaviour at
deployment, only to identify which path an answer depends on. The
patch-difference procedure of Section~\ref{sec:anatomy-diff} adds a small piece
of method here, turning per-patch logit-lens decodes into an exhaustive,
rankable list of changes rather than a chosen illustration.

\paragraph{Relation to our own prior work.} This paper extends a workshop paper
by the present authors \citep{anon-workshop}, included in the supplementary
material for reviewers. We state the overlap plainly. That paper established the
question-first paradox on NaturalBench, POPE and Winoground, diagnosed it as a
read-out failure using the perception probe, the read-out probe and the
attention knockout, proposed question and image echoing as the repair, and
reported the reversed-echo control. Sections~\ref{sec:paradox}
and~\ref{sec:anatomy} recap that material in compressed form, with all numbers
reproduced on our own runs, because the new results are not interpretable
without it. Everything in Sections~\ref{sec:results} to~\ref{sec:ablations} is
new: fourteen further benchmarks including video and object detection, the
reversal of the effect under a detection objective, the echo-resolution sweep
and its exact token accounting, the automated patch-difference procedure, and
the prompt-optimizer comparison. The earlier paper named the second image's
token cost as an open limitation and did not measure it; measuring it, and
showing how little of the echo is needed, is the core of this one.
